\documentclass[10pt,a4paper]{article}

% --- PAQUETES ---
\usepackage[utf8]{inputenc}
\usepackage[spanish, es-tabla]{babel}
\usepackage[margin=2.8cm, top=3cm, bottom=3cm, headheight=24pt]{geometry}
\usepackage[table]{xcolor}
\usepackage{amsmath, amssymb}
\usepackage{tabularx}
\usepackage{booktabs}
\usepackage{tcolorbox}
\tcbuselibrary{skins}
\usepackage{fancyhdr}
\usepackage{titlesec}
\usepackage{enumitem}
\usepackage{float}
\usepackage{lastpage}
\usepackage{graphicx}
\usepackage{caption}
\usepackage{listings}
\usepackage{hyperref}
\hbadness=10000

% --- CONFIGURACIÓN DE COLORES NEUTROS ---
\definecolor{PrimaryBlue}{RGB}{20, 60, 110}      % Azul oscuro corporativo neutro
\definecolor{LightBlue}{RGB}{240, 245, 252}      % Azul claro para tablas
\definecolor{WarningRed}{RGB}{170, 45, 45}       % Rojo oscuro para advertencias
\definecolor{LightRed}{RGB}{253, 240, 240}       % Rojo claro para fondo de nota
\definecolor{InfoGold}{RGB}{170, 130, 30}        % Dorado/Amarillo oscuro
\definecolor{LightGold}{RGB}{255, 253, 240}      % Amarillo claro para notas de formato

\definecolor{CodeBack}{RGB}{247, 247, 249}
\definecolor{CodeComment}{RGB}{90, 115, 90}
\definecolor{CodeKeyword}{RGB}{20, 60, 110}
\definecolor{CodeString}{RGB}{170, 45, 45}

% --- CONFIGURACIÓN DE CÓDIGO ---
\renewcommand{\lstlistingname}{Código}
\renewcommand{\lstlistlistingname}{Índice de códigos}

\lstdefinestyle{CustomStyle}{
    language=R,
    basicstyle=\ttfamily\footnotesize,
    keywordstyle=\color{CodeKeyword}\bfseries,
    commentstyle=\color{CodeComment}\itshape,
    stringstyle=\color{CodeString},
    backgroundcolor=\color{CodeBack},
    frame=single,
    rulecolor=\color{PrimaryBlue},
    breaklines=true,
    showstringspaces=false,
    numbers=left,
    numberstyle=\tiny\color{gray},
    tabsize=2,
    captionpos=b
}

% --- CONFIGURACIÓN DE ENCABEZADOS Y PIES DE PÁGINA ---
\pagestyle{fancy}
\fancyhf{}
\fancyhead[L]{\textcolor{PrimaryBlue}{\textbf{[Nombre de la Institución]}}}
\fancyhead[R]{\textcolor{PrimaryBlue}{[Documento / Informe Técnico]}}
\cfoot{\small Página \thepage\ de \pageref{LastPage}}
\renewcommand{\headrulewidth}{1pt}
\renewcommand{\headrule}{\hbox to\headwidth{\color{PrimaryBlue}\leaders\hrule height \headrulewidth\hfill}}

% --- CONFIGURACIÓN DE SECCIONES ---
\titleformat{\section}{\color{PrimaryBlue}\Large\bfseries}{\thesection.}{0.5em}{}
\titleformat{\subsection}{\color{PrimaryBlue}\large\bfseries}{\thesubsection.}{0.5em}{}

% --- CONFIGURACIÓN DE LISTAS ---
\setlist[itemize]{label=\textcolor{black}{\small$\blacksquare$}}
\setlist[enumerate]{leftmargin=*}

% --- INICIO DEL DOCUMENTO ---
\begin{document}
\hypersetup{pageanchor=false}

% --- PÁGINA 1: PORTADA ---
\thispagestyle{empty} 

% Banner superior
\noindent
\begin{tcolorbox}[colback=PrimaryBlue, colframe=PrimaryBlue, sharp corners, boxsep=3mm, halign=center]
    \color{white}
    \Large\textbf{[Nombre de la Institución / Facultad]} \\[1.5mm]
    \small [Departamento / Escuela Académica o Área Corporativa] \\
    \small [Asignatura / Proyecto de Referencia]
\end{tcolorbox}

\vspace{1.5cm}

% Título Central
\begin{center}
    \textcolor{PrimaryBlue}{\Huge \textsc{[Título del Informe / Proyecto]}} \\[0.5cm]
    \Large \textbf{[Subtítulo o Tema Principal del Análisis]} \\[0.2cm]
    \normalsize [Descripción resumida del contenido del documento] \\[1.5cm]
    
    \noindent\makebox[\textwidth]{\textcolor{PrimaryBlue}{\rule{0.6\textwidth}{1.5pt}}}
\end{center}

\vspace{2.5cm}

% Identificación
\begin{flushright}
    \begin{minipage}{0.5\textwidth}
        \normalsize
        \textbf{Autor(es):} \\
        [Nombre del Autor 1] \\
        [Nombre del Autor 2] \\[0.4cm]
        \textbf{Docente / Supervisor:} \\
        [Nombre del Docente / Evaluador] \\[0.4cm]
        \textbf{Fecha:} \\
        \today
    \end{minipage}
\end{flushright}

\clearpage
\hypersetup{pageanchor=true}

% --- PÁGINA 2: ÍNDICES ---
\tableofcontents
\newpage

% --- CONTENIDO ---
\section{Introducción}
Este documento sirve como plantilla formal en \LaTeX\ para informes técnicos, científicos o académicos. En esta sección se presenta la justificación de la investigación, el planteamiento del problema y la definición de objetivos.

El objetivo principal de este informe se establece a continuación:
\begin{center}
    \textit{\textbf{[Definir aquí el objetivo principal del análisis en formato destacado y negrita]}}
\end{center}

\section{Metodología y Datos}
En esta sección se describe la muestra analizada, la naturaleza de las variables y el tratamiento metodológico aplicado.

\subsection{Estructura de la Base de Datos}
A continuación se presenta un cuadro descriptivo de las variables analizadas como ejemplo de maquetación de tablas usando el entorno \texttt{tabularx}:

\begin{table}[H]
\centering
\small
\begin{tabularx}{\textwidth}{l c X}
\toprule[1.5pt]
\rowcolor{PrimaryBlue} \textcolor{white}{\textbf{Variable}} & \textcolor{white}{\textbf{Tipo}} & \textcolor{white}{\textbf{Descripción y Codificación}} \\
\midrule
\texttt{Variable\_1} & Numérica & Variable de respuesta continua que representa [Unidad de Medida]. \\
\rowcolor{LightBlue}
\texttt{Variable\_2} & Categórica & Variable binaria de clasificación (0 = Control, 1 = Tratamiento). \\
\texttt{Variable\_3} & Numérica & Explicativa o covariable independiente (edad, meses, etc.). \\
\rowcolor{LightBlue}
\texttt{Variable\_4} & Categórica & Variable multinivel representando estados o categorías. \\
\bottomrule[1.5pt]
\end{tabularx}
\caption{Diccionario de variables e identificación del tipo de dato.}
\label{tab:variables}
\end{table}

\subsection{Modelamiento Matemático}
Es posible incorporar ecuaciones matemáticas avanzadas. Por ejemplo, un modelo matemático lineal se define mediante la ecuación:
\begin{equation}
    Y_i = \beta_0 + \beta_1 X_{1i} + \beta_2 X_{2i} + \varepsilon_i
\end{equation}
Donde $Y_i$ es la respuesta, $X_{ji}$ son los predictores y $\varepsilon_i$ representa el término de error aleatorio.

\section{Resultados y Diagnóstico}
En esta sección se exponen los hallazgos principales y la verificación de supuestos estadísticos correspondientes.

\subsection{Visualización de Datos}
Las figuras deben ser centradas e incluir su respectivo identificador y leyenda explicativa:

\begin{figure}[H]
    \centering
    % \includegraphics[width=0.6\textwidth]{imagenes/grafico_ejemplo.png}
    \begin{tcolorbox}[colback=white, colframe=PrimaryBlue, width=0.6\textwidth, halign=center, valign=center, height=4cm]
        \textcolor{PrimaryBlue}{\textbf{[Marcador de Posición para Gráfico / Imagen]}} \\
        (Ej: Boxplot de distribución o Matriz de correlación)
    \end{tcolorbox}
    \caption{Ejemplo de gráfico o visualización de datos.}
    \label{fig:grafico1}
\end{figure}

\section{Conclusiones}
Resumen de los principales resultados, aportes y posibles líneas de trabajo futuras.

\newpage
\section{Anexo: Código de Implementación}
Sección destinada a la presentación del código fuente del análisis (R, Python, etc.) debidamente comentado y estructurado:

\begin{lstlisting}[style=CustomStyle, caption={Script de ejemplo para análisis descriptivo genérico.}]
# Carga de librerías básicas
library(ggplot2)

# Análisis descriptivo exploratorio
data(iris)
summary(iris)

# Gráfico de ejemplo
ggplot(iris, aes(x = Sepal.Length, y = Sepal.Width, color = Species)) +
  geom_point(size = 3) +
  theme_minimal() +
  labs(title = "Análisis de Longitud y Ancho de Sépalo",
       x = "Longitud de Sépalo",
       y = "Ancho de Sépalo")
\end{lstlisting}

\end{document}
